\chapter{Introduction}
\label{chap:introduction}

\section{Motivation}
\label{sec:motivation}

The pace at which biological sequence data accumulates has long outstripped our
capacity to characterise proteins experimentally. As of early 2026, \gls{uniprot}
contains more than 250 million sequences, yet fewer than 0.02\,\% carry
experimentally validated functional annotations~\cite{uniprot2023}. The gap is not
merely quantitative: the proteins that remain unannotated are disproportionately
found in under-studied organisms, precisely where novel biology, and potential
therapeutic targets, are most likely to reside.

Computational function prediction attempts to bridge this gap by transferring
knowledge from characterised proteins to unknown ones. Classical approaches based on
sequence similarity (BLAST, profile \glspl{hmm}) remain widely used, but their
accuracy degrades for distantly related sequences. Over the last five years, however,
\glspl{plm} trained on hundreds of millions of sequences have produced dense vector
representations that capture structural and functional signal well beyond what
classical alignment can detect, and a generation of structure-conditioned models
(ESM-3, ProstT5) has begun to integrate sequence and structure into a single
representation. Each new PLM unlocks a new operating point in the accuracy/throughput
trade-off, but also introduces a new dependency, a new tokeniser, a new GPU memory
footprint, and a new failure mode. Comparing PLMs at scale, on a fair benchmark, with
reproducible numbers, is itself a research engineering problem.

PROTEA was designed to address both faces of the problem at once: as a method, it
operationalises embedding-based annotation transfer with selective learning-to-rank
re-finement; as a platform, it provides the extensible foundation that lets new
backends, annotation sources, and ranking methods be plugged in without ad-hoc
glue. The platform is deliberately built around the assumption that the canonical
pipeline of today will be obsolete tomorrow. What must remain stable is the contract
between layers, the trace of every run, and the pathway from a FASTA file to a
ranked, interpretable Gene Ontology prediction.

\section{Research Questions}
\label{sec:research-questions}

This thesis addresses the following questions:

\begin{enumerate}[label=\textbf{RQ\arabic*.}]
  \item Can embedding-based \gls{knn} annotation transfer match or exceed
        alignment-based methods for \gls{go} term prediction, particularly at low
        sequence identity, and how does the answer change as new PLMs replace
        older ones?
  \item Do pairwise alignment statistics, taxonomic distance, ontology-graph
        ancestry features, and per-PLM PCA projections provide complementary signal
        that improves the ranking of candidate annotations once the embedding
        baseline is already strong?
  \item How can a reproducible, scalable annotation pipeline be designed so that
        researchers can re-run analyses as ontologies and databases are updated, and
        so that the pipeline itself can be deployed in development, cloud, on-premise
        HPC, and airgapped environments without code branching?
  \item How can the journey of a research platform, with its dead ends and pivots,
        be captured operationally so that the canonical pipeline is justified by an
        auditable trace rather than by post-hoc rationalisation?
\end{enumerate}

\section{Contributions}
\label{sec:contributions}

The main contributions of this work are:

\begin{itemize}
  \item \textbf{PROTEA}: an open-source distributed platform implementing the full
        annotation pipeline, from FASTA upload to structured \gls{go} predictions
        with optional feature engineering and selective learning-to-rank re-ranking.
  \item A \textbf{plugin architecture} based on Python \texttt{entry\_points},
        organised across seven code repositories, that lets new annotation sources,
        runners and PLM backends be added without modifying the core platform.
  \item A \textbf{versioned data model} for ontology snapshots, annotation sets,
        embedding configurations, datasets, re-ranker models and experiment runs,
        with provenance complete enough to exactly replay any past prediction.
  \item A \textbf{benchmark of eight PLM backends} (ESM-2 in three sizes, ESM-C,
        ProstT5, ProtT5, Ankh, ESM-3 component encoder) on a common evaluation
        surface, with per-aspect, per-category and per-backend numbers.
  \item A \textbf{feature engineering layer} that augments embedding similarity with
        Needleman--Wunsch / Smith--Waterman alignment scores,
        NCBI-taxonomy-based distance metrics, GO-graph ancestry embeddings, and
        per-PLM PCA projections, organised as a registry of approximately 52
        features.
  \item A \textbf{selective learning-to-rank re-ranker} trained per category and per
        aspect, that falls back to the embedding-only baseline whenever re-ranking
        would degrade performance for the current cell.
  \item A \textbf{multi-target deployment story} (development \texttt{docker
        compose}, cloud Helm/Compose, HPC Apptainer, airgapped bundle) covering the
        same canonical pipeline without code branching.
  \item An \textbf{operational narrative model} (\texttt{Job.findings},
        \texttt{ExperimentRun}, comments) that captures the why of every run and
        provides the audit trace from which the evaluation chapter is distilled.
  \item An \textbf{empirical evaluation} comparing PROTEA predictions against
        baseline methods on held-out \gls{goa} splits following the CAFA temporal
        protocol, complemented by three submissions to the LAFA public benchmark.
\end{itemize}

\section{Thesis Structure}
\label{sec:structure}

The remainder of this document is organised as follows.
\Cref{chap:background} introduces the biological concepts underpinning the work,
including the modern landscape of protein language models and their structural
extensions.
\Cref{chap:related} surveys existing computational approaches to protein function
prediction up to 2025, with attention to the most recent embedding-based and
graph-based methods.
\Cref{chap:design} presents the architecture of PROTEA, its plugin-based
extensibility model, and the operational narrative layer.
\Cref{chap:implementation} describes key implementation decisions across the seven
repositories, including the configuration system, the artifact-store boundary, the
queue topology, and the deployment containers.
\Cref{chap:evaluation} reports the experimental results on the eight-PLM benchmark,
the impact of feature families through ablation, and the journey that shaped the
canonical configuration.
\Cref{chap:conclusion} summarises the findings, discusses limitations, and outlines
future directions including the LAFA submissions and the post-defense agenda.
